\section{Introduction}

The forced internment of over 120,000 Japanese Americans during World War II 
stands as a stark episode of racial injustice and economic loss. 
The policy displaced individuals and families 
from their homes, businesses, and farms on the West Coast;
a forced displacement that fractured communities and interrupted lives. 
Conventional understanding rightly frames this event as a traumatic scar, 
with research documenting its detrimental effects 
on property, wealth, and psychological well-being 
\citep{commissiononwartimerelocationPersonalJusticeDenied1983}.

However, an emerging body of economic literature
presents a more complex and surprising picture. 
Despite losses of property, job opportunities, and livelihoods,
some evidence suggests that for a subset of the internee population
this forced displacement paradoxically led to upward economic mobility 
in the decades following the war
\citep{arellano-boverDisplacementDiversityMobility2022,shoagCausalEffectPlace2016}.
This finding presents a central puzzle: 
How could a policy of forced removal and incarceration 
have created economic \textit{opportunity}? 
In this paper I examine the mechanisms through which 
internment reshaped the labor market outcomes of Japanese Americans 
and, crucially, for whom these effects were most pronounced.

I argue that, at least for some internees,
internment acted as a large-scale, "move to opportunity" experiment.
Forced relocation may have generated economic gains through two primary channels.
First, it allowed internees to escape discrimination on the West Coast,
where the anti-Japanese sentiment that originally contributed to internment
may have held back their potential post-war earnings. 
Second, the quasi-random placement of internees in camps across the interior U.S.
inadvertently moved them closer to burgeoning post-war economies, 
such as those in the Midwest and Chicago,
creating new labor market prospects that would have otherwise been inaccessible
\citep{austinEastwardPioneersJapanese2007}.

To test this hypothesis, 
I use newly available linked census data from IPUMS 
to track individuals from 1940 to 1950.
I combine these census records with detailed War Relocation Authority (WRA) 
records on internee origins and camp assignments. 
My empirical strategy leverages the intensity of internment at the county level 
to identify its causal impact. 
By comparing the outcomes of Japanese Americans 
from counties with high internment rates to a control group of Chinese Americans
and non-West Coast Japanese\footnote{Those living in the US mainland outside of Arizona, California, Oregon, or Washington during 1940.}
who were not targeted by the policy,
I can estimate the effect of forced displacement on post-war earnings
and migration.

My analysis yields three key findings. 
First,
greater internment exposure predicts higher 1950 earnings,
consistent with a place-based opportunity effect.
Second, this effect was not uniform across the population; 
it was driven almost entirely by women. 
Japanese American women who were likely internees 
saw substantially larger wage gains than their non-interned peers.
This result highlights a critical intersection between forced migration 
and gender dynamics. 
Third, I find that these gains are linked to migration; 
internees were more likely to move to
more densely populated counties with higher median incomes.

This paper contributes in three main ways.
First, it adds supporting evidence to the long-term economic consequences 
of Japanese internment 
and clarifies potential mechanisms behind
the surprising positive effects for labor market outcomes
found in the economic literature.
Second, this study leverages newly linked historical panel data
to study large-scale displacements of an ethnic minority group.
Third, it contributes to the broader place-effects literature 
by showing how a large, exogenous shock to an individual's location 
can alter their economic trajectory, 
particularly for historically marginalized groups. 
The findings shed light on the complex legacy of internment 
and the different ways in which marginalized people can adapt
in spite of discrimination and adversity.
Importantly, the findings of this paper should not be seen
as diminishing the moral and material harms of internment overall.
My dataset does not measure property losses,
or emotional trauma from unseating an entire population from their homes.
Instead, this paper shows that despite all of these losses,
internees were able to adapt in a relatively short period of time
to the new locations and opportunities
in which they found themselves after their release.

\subsection{\label{sec:history}Historical Background}
\href{https://www.archives.gov/milestone-documents/executive-order-9066}{Executive Order 9066} 
was signed on February 19, 1942 and granted the Secretary of War the authority
to designate military areas from which civilians could be prevented access. This
order was used to target all persons of Japanese descent living in the so-called
`Exclusion Zone' which included all of California, and large areas of
Washington, Oregon, and Arizona shown in blue in Figure \ref{fig:internmentmap}.
At first, the Secretary of War Henry Stimson encouraged some families 
to relocate voluntarily from the West Coast,
but starting at the end of March 1942, 
the Wartime Civil Control Administration (WCCA) and War Relocation Authority (WRA)
were established to remove all remaining Japanese Americans to temporary
Assembly Centers which were makeshift shelters in close proximity to the
population centers where these people were living. As the more permanent WRA
Relocation Centers were completed in May 1942, all Japanese Americans left
on the West Coast were moved from Assembly Centers to Relocation Centers by
train. This process lasted until November of 1942
\citep{commissiononwartimerelocationPersonalJusticeDenied1983, denshoencyclopediacontributorsResettlementDenshoEncyclopedia2020}
.

As seen in Figure \ref{fig:1940JAmap}, over 90\% of the 1940 Japanese American
population lived in areas which would become subject to relocation orders.
This population was overwhelmingly concentrated along the West Coast, 
particularly in California,
which was home to large, well-established Japanese American communities,
many of whom were engaged in farming and small businesses.
The geographic concentration of this population 
in economically vital areas near major cities and military installations 
heightened the perceived security threat,
prompting their forced relocation.

Importantly for the later identifying assumptions in my empirical methods,
there is historical anecdotal evidence to suggest that these
voluntary migrants may have been an economically selected group.
Some Japanese living on the West Coast were permitted to leave the Exclusion Zone
if they had a sponsor living in another state who was able to prove that
they had a job or educational program in that state.
Some Japanese families were able to rely on business connections.
For example, in an interview with the 
\href{https://ddr.densho.org/media/ddr-manz-1/ddr-manz-1-34-transcript-2e0e58d5a9.htm}{Densho Digital Archive},
Art Imagire recalls his mother's decision to move their family to Reno, Nevada 
because one of her former clients offered her a job as a dress-maker there 
\citep{imagireArtImagireInterview2008}.
Imagire cites his mother's value for her children's education
as an important reason for why their family was one of the few which chose,
``to drop everything and move''.
Additionally, several thousand Nisei were attending college 
when Pearl Harbor was attacked.
Many Nisei students such as Sam Naito 
were sent home from West Coast schools like the University of Oregon
during their academic term \citep{naitoSamNaitoInterview2003}.
Some, like Naito, were able to transfer to schools outside of the exclusion
zone to continue their education.
In my empirical strategy, I address this possibility that
some Japanese may have chosen to voluntarily migrate
by controlling for their observable 1940 characteristics
and by comparing against West Coast Chinese in the control group.

After two years of WRA camp operations, the government began to debate how to
end the evacuation orders towards the end of 1944. One statement issued on
December 9 illustrates that one of the initial goals of interment was actually
to cause a permanent shift in the Japanese American population away from the
West Coast as part of a wide-scale program of `deghettoization':
\begin{quote}
  ``One of the major WRA aims, from the beginning, has been to encourage the widest possible dispersal of evacuees throughout the Nation, and this will continue as a prime objective during the final phase of the program.''
  \hfill \citetalias[p.~234]{commissiononwartimerelocationPersonalJusticeDenied1983}
\end{quote}

To the extent to which the relocation policies described above were successful,
this period of history provides an opportunity to examine
the way that contemporary or future displaced people
may potentially adapt to new labor markets and conditions.

\subsection{Motivating Literature}\label{motivation}

Several existing studies have tried to quantify the causal effect 
of Japanese internment on long-term earnings
\citep{chinLongRunLaborMarket2005}, 
childhood health
\citep{saavedraEarlyChildhoodConditions2013, 
grossmanIntergenerationalHealthEffects2025},
and educational attainment 
\citep{saavedraSchoolQualityEducational2015}. 

On the other hand, other studies have identified surprising ways 
in which internees did better than their counterparts.
\cite{shoagCausalEffectPlace2016} examine the persistent pattern of 
internees' tendency to settle closer to their assigned camp,
and as a result find that differential exposure to more affluent areas 
led to some internees experiencing more long-term upward mobility than others.
\cite{arellano-boverDisplacementDiversityMobility2022} finds that overall, 
interned Japanese men experienced 9-22\% higher incomes 
relative to similar Chinese and non-interned Japanese men.
He attributes this effect to a combination of information and skills transfers 
among those living in camps together, 
migration to new labor markets, and career switching
that would not have happened if internees had remained in their original homes 
on the West Coast.

These economic findings are reinforced by more qualitative evidence
voiced by the Nisei themselves
and in historical research on this period.
\cite{taylorLeavingConcentrationCamps1991}
points to the differences in the benefits of relocation between
those who left earlier (1942-1944) via work or education programs,
and those who left in 1944 and 1945 who tended to be older
and more resistant to the government's policies.
\cite{austinEastwardPioneersJapanese2007} 
highlights the experiences of several Nisei who resettled in cities such as
Chicago, Detroit and Minneapolis.
For those who faced discrimination on the West Coast
where anti-Japanese sentiment was more prevalent,
these new cities felt more socially welcoming
and economically promising.

Aside from the importance of understanding the losses to internees, Japanese
Internment can also serve as an interesting natural experiment through the
involuntary nature of the experiences of internees. 
A key question in previous economic literature is the extent to which 
geographic location influences economic outcomes, such as upward mobility or earnings.
It is difficult to provide an unbiased causal estimate because households usually
have a large degree of choice in where they live. Previous studies have used a
variety of empirical strategies to separate out this causal effect including
controlling for observable characteristics like education
\citep{glaeserCitiesSkills2001}, worker fixed-effects or
through leveraging other exogenous placements of households like the Moving to
Opportunity Experiment 
\citep{chettyEffectsExposureBetter2016}. 

Related to the overall racial discrimination faced by Japanese Americans,
Chinese Americans also faced immigration barriers
\citep{chenImpactSkillBasedImmigration2015},
and discrimination
\citep{chenInstitutionalDiscriminationAssimilation2024}.

Other research points to the importance of social networks for migration choices.
\cite{munshiSocialNetworksMigration2020}
expands the standard Roy model of migration choice
\citep{royThoughtsDistributionEarnings1951}
by adding the possibility for the number of individuals in a potential migrant's
social network living in a given destination to affect their utility
of living there themselves.

% *** Expand on gender differences at this time
% to explain heterogeneity findings
